% Paper 3, §9 — Test-time scaling beyond CoT: think-tokens, budget forcing, reasoning architectures
% Anchors: kb/notes/reasoning/test-time-compute.md §2.3 (sequential TTC)
%          kb/notes/reasoning/reasoning-training.md §2 (R1 four-stage protocol)
%          kb/excerpts/s1-2025.md (verbatim source)
% Cross-paper: Paper-1 §reasoning-architectures (architecture taxonomy),
%              Paper-2 §SFT / §post-training (training taxonomy).

\section{Test-time scaling beyond CoT: think-tokens, budget forcing, reasoning architectures}
\label{sec:budget-forcing}

\Cref{sec:test-time-compute,sec:tree-search} treat inference-time
compute as a \emph{sampling} or \emph{search} budget --- $N$ parallel
traces, $B$ branches, $N_\text{rollout}$ MCTS simulations. A third
family holds $N = 1$ but lengthens the single trace, spending compute
inside one \glsterm{autoregressive}{autoregressive} decode rather than across many. This is the
o1 / R1 / QwQ / Gemini-2.5-thinking
paradigm.\footnote{KB note: \texttt{kb/notes/reasoning/test-time-compute.md\#2.3}.}
\citet{s1-2025} establish the minimum recipe for it: $1{,}000$
curated long-\glsterm{chain-of-thought}{CoT} exemplars for SFT, plus a \textbf{\glsterm{budget-forcing}{budget forcing}} rule
that appends the literal token \texttt{Wait} to extend thinking and
appends an end-of-thinking delimiter to truncate it
\citep[\S 3.1]{s1-2025}.\footnote{KB excerpt:
\texttt{kb/excerpts/s1-2025.md\#sec-3-1-budget}.} The exercise isolates
\emph{which} of the moving parts in a reasoning-trained model
(\cref{sec:r1-family}) is doing the test-time scaling work, and in
doing so locates a \textbf{reasoning-architecture} axis that is
separate from the base-architecture axis of Paper~1: the same
\glsterm{transformer}{Transformer}, but a training-and-inference protocol that designates
\texttt{<think>...</think>} as a structurally distinct phase from the
answer phase. This section formalises the thinking-token budget $L$ as
the operative compute knob, transcribes the s1 mechanism, draws the
training-time vs inference-time line that separates s1 from R1 cold
start, and forwards to \cref{sec:search-vs-rl}'s fixed-compute
trade-off.

\subsection{Thinking-token budget $L$}
\label{sec:budget-knob}

Let $x$ be a problem and $z = (z_1, \ldots, z_L)$ a thinking trace
emitted before the answer phase. \emph{Sequential} \glsterm{test-time-compute}{test-time compute}
holds $N = 1$ but treats $L$ as the explicit compute knob:
\begin{equation}
  z \sim p_\theta\bigl(\cdot \mid x;\; \text{length budget } L\bigr),
  \qquad
  y \sim p_\theta\bigl(\cdot \mid x,\; z\bigr).
  \label{eq:sequential-ttc}
\end{equation}
The compute cost grows in $L$ along two terms: a per-token decode that
is itself $O(L)$ in the prefix-attention pass, so total cost is
$\Theta(L^2)$ in the trace-length-bound regime where the \glsterm{kv-cache}{KV cache}
dominates.\footnote{KB note:
\texttt{kb/notes/reasoning/test-time-compute.md\#2.3}.} At inference
time the trace is bracketed by special tokens
\texttt{<think>}~$\ldots$~\texttt{</think>}, after which the model
emits the answer; the brackets are part of the chat template and the
\texttt{</think>} delimiter is the model's signal that thinking is
complete.

This is the third \glsterm{test-time-compute}{TTC} family of \cref{sec:test-time-compute}: not
$N$ samples (sampling, \cref{sec:self-consistency}), not branching
(search, \cref{sec:tree-search}), but length. The two control surfaces
exposed in production are:
\begin{enumerate}
  \item \textbf{API-level thinking budgets.} Gemini~2.5 takes a
        \texttt{thinkingBudget} integer; OpenAI's o-series exposes
        \texttt{reasoning\_effort} levels; Anthropic exposes a
        \texttt{thinking} block with token caps.\footnote{KB note:
        \texttt{kb/notes/reasoning/test-time-compute.md\#2.3} §1.}
        The user sets a target $L$; the serving stack honours it.
  \item \textbf{\glsterm{budget-forcing}{Budget forcing}.} A purely test-time edit to the
        decoding loop, requiring no API surface and no retraining
        \citep[\S 3.1]{s1-2025}. We turn to it next.
\end{enumerate}

\subsection{The s1 budget-forcing rule}
\label{sec:s1-rule}

\citet[\S 3.1]{s1-2025} define \glsterm{budget-forcing}{budget forcing} as two single-line
edits to \glsterm{autoregressive}{autoregressive} decoding, parameterised by a minimum budget
$L_\text{min}$ and a maximum budget $L_\text{max}$. Let
$\ell$ denote the current trace length and let $\langle/\text{think}\rangle$
denote the end-of-thinking token. Then:
\begin{align}
  \text{(min)} \quad &\text{if } \ell < L_\text{min} \text{ and } \arg\max_v\, p_\theta(v \mid x, z_{1:\ell}) = \langle/\text{think}\rangle:\nonumber \\
  &\quad p_\theta(\langle/\text{think}\rangle \mid \cdot) \leftarrow -\infty,\;\; z_{\ell+1} \leftarrow \texttt{Wait}, \label{eq:bf-min} \\
  \text{(max)} \quad &\text{if } \ell \ge L_\text{max}:\quad z_{\ell+1} \leftarrow \langle/\text{think}\rangle. \label{eq:bf-max}
\end{align}
The min-rule (\cref{eq:bf-min}) suppresses the end-of-thinking token's
probability and appends the literal string \texttt{Wait} to the trace
when the model tries to stop too early; the max-rule
(\cref{eq:bf-max}) hard-injects the end-of-thinking delimiter when the
trace runs too long. The \texttt{Wait} insertion is the load-bearing
piece: rather than a benign continuation cue, it is a \emph{prompt} to
re-examine the work just produced
\citep[\S 3.1]{s1-2025}.\footnote{KB excerpt:
\texttt{kb/excerpts/s1-2025.md\#sec-3-1-budget}.} The paper's verbatim
description: ``This can lead the model to double-check its answer,
often fixing incorrect reasoning steps''
\citep[abstract]{s1-2025}.\footnote{KB excerpt:
\texttt{kb/excerpts/s1-2025.md\#abstract}.}

The recipe is then three lines:
\begin{enumerate}
  \item Curate $1{,}000$ (problem, long-\glsterm{chain-of-thought}{CoT} trace, answer) examples
        filtered by quality, difficulty, and diversity
        \citep[\S 2.2]{s1-2025}.\footnote{KB excerpt:
        \texttt{kb/excerpts/s1-2025.md\#sec-2-2-curation}.}
  \item SFT a Qwen2.5-32B-Instruct base on these $1{,}000$ traces
        ($\sim 7$ H100-GPU-hours).
  \item Apply \glsterm{budget-forcing}{budget forcing} at test time per
        \cref{eq:bf-min,eq:bf-max}.
\end{enumerate}
No \glsterm{reward-model}{reward model}. No RL. No search tree. The budget-forcing rule itself
is $\sim 10$ lines of decoder code.

\begin{intuition}
\glsterm{budget-forcing}{Budget forcing} lives on Snell's test-time-compute axis
(\cref{sec:test-time-compute}) but exercises it from the
\emph{sequential} side: where Snell varies $N$ at fixed $L$, s1 varies
$L$ at $N = 1$. Operationally, \texttt{Wait} tokens are the
sequential analogue of resampling --- each \texttt{Wait} forces the
model to revisit its own partial trace and emit a refinement, much as
each new sample in \glsterm{self-consistency-2}{self-consistency} (\cref{sec:self-consistency})
draws a fresh attempt at the problem. Returning to math: the compute
expended per problem is $C \propto L$ for sequential and $C \propto N$
for sampling, and the two axes can be swept independently. s1's
ablation \citep[\S 5.1]{s1-2025}\footnote{KB excerpt:
\texttt{kb/excerpts/s1-2025.md\#sec-5-1-ablation}.} shows the
sequential axis is non-trivial in isolation: AIME~2024 lifts from
$50.0\%$ (SFT only, no \glsterm{budget-forcing}{budget forcing}) to $56.7\%$ (SFT plus budget
forcing) on the same checkpoint. The lift is what \glsterm{budget-forcing}{budget forcing}
\emph{adds} on top of SFT --- the sequential \glsterm{test-time-compute}{TTC} dial, exposed as a
post-hoc decoder edit.
\end{intuition}

\subsection{Logarithmic-ish returns to $L$}
\label{sec:log-returns}

The s1 paper sweeps $L$ through \glsterm{budget-forcing}{budget forcing} and reports
extrapolation \emph{beyond} the SFT trace lengths
\citep[abstract, \S 4.2]{s1-2025}: pushing the minimum-\glsterm{thinking-budget-2}{thinking budget}
upward continues to lift accuracy past the lengths the model saw in
training. The concrete numbers from
Table~1 \citep[\S 4.2]{s1-2025}:\footnote{KB excerpt:
\texttt{kb/excerpts/s1-2025.md\#sec-4-2-table}.}

\begin{table}[h]
\centering
\small
\begin{tabular}{lrrr}
\toprule
Model & AIME24 & MATH500 & \glsterm{gpqa-diamond-2}{GPQA Diamond} \\
\midrule
s1-32B (with \glsterm{budget-forcing}{budget forcing}) & $56.7\%$ & $93.0\%$ & $59.6\%$ \\
o1-preview                    & $44.6\%$ & $85.5\%$ & $73.3\%$ \\
Qwen2.5-32B-Instruct (base)   & $26.7\%$ & $84.0\%$ & $49.0\%$ \\
\bottomrule
\end{tabular}
\caption{s1-32B vs.\ o1-preview and the underlying Qwen2.5-32B-Instruct
base \citep[Table~1, \S 4.2]{s1-2025}.\protect\footnotemark{}
On competition-math benchmarks (AIME24, MATH500), $1{,}000$
SFT examples plus budget forcing exceeds o1-preview by
$+12.1$ and $+7.5$ percentage points respectively; on
GPQA Diamond, o1-preview retains a substantial lead, indicating that
the recipe transfers to math-style reasoning more readily than to
broader scientific QA.}
\label{tab:s1-headline}
\end{table}
\footnotetext{KB excerpt: \texttt{kb/excerpts/s1-2025.md\#sec-4-2-table}.}

The shape of accuracy as a function of $L$ on the s1 curves is
\emph{monotone-concave} in $\log L$ over the range tested
\citep[\S 4.2]{s1-2025}; no closed-form \glsterm{scaling-law}{scaling law} is reported. This
is consistent with the broader Snell-style framing of
\cref{sec:test-time-compute} that posits inference-time compute as a
scaling axis comparable to training compute, but with the functional
form left empirical \citep[\S 1]{snell2024}. The headline ablation
\citep[\S 5.1]{s1-2025} additionally controls for the
\emph{interaction} between SFT data scale and the budget-forcing
rule: training on the full $59{,}000$-example pool reaches $53.3\%$
AIME~2024 in $\sim 394$ H100-GPU-hours, while the curated $1{,}000$
plus \glsterm{budget-forcing}{budget forcing} reaches $56.7\%$ in
$\sim 7$ H100-GPU-hours.\footnote{KB excerpt:
\texttt{kb/excerpts/s1-2025.md\#sec-5-1-ablation}.} Curation does most
of the SFT-side work; \glsterm{budget-forcing}{budget forcing} extracts the rest at inference.

\begin{contradiction}
The s1 thesis --- $1{,}000$ examples plus \glsterm{budget-forcing}{budget forcing}
\emph{suffices} for o1-preview-level competition math --- is in
explicit tension with the DeepSeek-R1 framing
(\cref{sec:r1-family}) that \emph{RL on top of SFT} is required for
the strongest performance \citep[\S 2.2.4]{deepseek-r1}.\footnote{KB
note: \texttt{kb/notes/reasoning/test-time-compute.md\#3} flags this
contradiction explicitly.} The clearest reconciliation in the 2025--26
literature: SFT on curated traces gets the model into the right
regime cheaply; RL refines within that regime but cannot create the
regime from scratch on a base that lacks the latent capability. s1
operates within the regime; R1 grows the regime via outcome-reward
RL. Both can be right because they target different points on the
training-compute axis.
\end{contradiction}

\subsection{Inference-time discipline vs.\ training-time discipline}
\label{sec:bf-vs-cold-start}

\glsterm{budget-forcing}{Budget forcing} and R1's cold-start SFT (\cref{sec:r1-family}) are
\emph{both} mechanisms that shape the same \texttt{<think>} channel,
but at different stages of the pipeline. The distinction is
load-bearing for cross-paper bookkeeping:

\begin{itemize}
  \item \textbf{\glsterm{budget-forcing}{Budget forcing} is inference-time.}
        \cref{eq:bf-min,eq:bf-max} change the decoding loop, not the
        weights. The model is whatever the SFT step left behind; the
        rule is a post-hoc dial on $L$ at serve time. This is the same
        category as Gemini's \texttt{thinkingBudget} or OpenAI's
        \texttt{reasoning\_effort} --- a knob exposed to the user or
        the serving stack, not baked into $\theta$.

  \item \textbf{\glsterm{cold-start-sft-2}{Cold-start SFT} is training-time.} R1's Stage~1
        \citep[\S 2.3]{deepseek-r1} fine-tunes the base model on
        long-\glsterm{chain-of-thought}{CoT} exemplars before any RL, ensuring readable traces and
        a starting point that matches the chat template's
        \texttt{<think>}~$\ldots$~\texttt{</think>}
        format.\footnote{KB note:
        \texttt{kb/notes/reasoning/reasoning-training.md\#2} (Stage~1).}
        It changes $\theta$. Whatever distribution-shape the base
        model has after Stage~1 is what every later mechanism ---
        Stage~2 \glsterm{grpo-2}{GRPO}, Stage~3 \glsterm{rejection-sampling-sft}{rejection-sampling SFT}, and any
        inference-time budget rule --- operates on top of.
\end{itemize}

The two compose. s1 exhibits the inference-time leg in isolation by
running on a Qwen2.5-32B-Instruct base whose pretraining and
instruction tuning supply the latent reasoning capability
\citep[\S 5.1]{s1-2025}. R1 exhibits the training-time leg in
isolation through \glsterm{r1-zero}{R1-Zero}, where \glsterm{grpo-2}{GRPO} from the base model produces the
trace-lengthening behaviour without any cold-start SFT
\citep[\S 2.2.4]{deepseek-r1}. The frontier-2026 picture
(\cref{sec:frontier-snapshot}) is models that ship \emph{both} ---
training-time cold-start to seed the regime, then inference-time
budget exposure as a productionised dial.

\begin{intuition}
The architecture taxonomy of Paper~1 (Paper~1 §reasoning-architectures)
treats reasoning as an architectural element by asking what changes in
the \glsterm{transformer}{Transformer} when it must emit and consume long traces:
long-context attention, KV-cache reuse, position-encoding behaviour at
extrapolated lengths. The training taxonomy of Paper~2 (Paper~2
§SFT and §post-training) treats it as a training-pipeline element:
which datasets, which losses, which RL stages produce the trace-
emitting policy. \emph{Reasoning-architecture}, in the sense
introduced by s1 and R1 jointly, is a third axis: the
training-and-inference \emph{protocol} that designates
\texttt{<think>}~$\ldots$~\texttt{</think>} as a structurally distinct
phase, separates thinking-token budgets from answer-token budgets, and
exposes $L$ as a serve-time control surface. Returning to math: the
underlying model is the same $\theta \in \R^{|\theta|}$ as in
Paper~1; the protocol is a constraint on the input-output trajectory
$(x, z, y)$ that turns $L = |z|$ into an additive compute axis on top
of the architectural and training axes.
\end{intuition}

\subsection{Reasoning-architecture as a third axis}
\label{sec:reasoning-arch-axis}

The cross-paper claim worth stating sharply: \glsterm{budget-forcing}{budget forcing} is the
\emph{simplest possible} demonstration that reasoning-architecture is a
distinct degree of freedom from base architecture and from training
recipe. The same Qwen2.5-32B-Instruct base that
\cref{tab:s1-headline} reports at $26.7\%$ AIME~2024 zero-shot lifts
to $50.0\%$ under SFT-on-1K-traces and to $56.7\%$ under SFT plus
\glsterm{budget-forcing}{budget forcing} \citep[\S 4.2, \S 5.1]{s1-2025}. The base architecture
did not change. The training-data scale changed by a factor of
$\sim 30$ relative to the 59K-pool ablation but only marginally
relative to instruction tuning. The reasoning-architecture --- the
chat-template separation of thinking from answering, plus the rule
that lets a serve-time dial control how much thinking happens ---
absorbs the rest. The rest is $\sim 30$ percentage points of
AIME~2024 pass@1.

This locates two cross-paper bookkeeping items:
\begin{enumerate}
  \item In the architecture taxonomy of Paper~1
        (Paper~1 §reasoning-architectures), reasoning shows up as
        long-context attention quality and KV-cache costs --- the
        \emph{enablers} of long-trace inference. The
        \texttt{<think>} bracket itself is a chat-template artifact,
        not an architectural primitive; it lives at the
        protocol layer.
  \item In the training taxonomy of Paper~2 (\hyperref[P2-sec:sft]{Paper~2 §SFT},
        §post-training), reasoning shows up as the curated long-\glsterm{chain-of-thought}{CoT}
        SFT stage and the \glsterm{rlvr-3}{RLVR} stage that follows
        (\cref{sec:rlvr,sec:r1-family}). The \emph{format} of the
        traces (delimiters, line-by-line structure, inline
        verification) is part of what SFT installs; what RL does is
        push toward longer and more verifier-aligned versions of that
        format \citep[\S 2.2.2]{deepseek-r1}.
\end{enumerate}
The third axis --- the reasoning-architecture / serve-time-protocol
layer --- is where \glsterm{budget-forcing}{budget forcing}, \texttt{thinkingBudget} APIs,
\texttt{reasoning\_effort} levels, and Anthropic's \texttt{thinking}
blocks all live. None of them changes the model weights; all of them
modify what $L$ the model is allowed (or required) to spend.

\subsection{Forward link}
\label{sec:bf-forward}

\glsterm{budget-forcing}{Budget forcing} exposes $L$ as a clean inference-time knob, but it
leaves the central operational question of the reasoning triangle
unanswered: at fixed total compute, when does spending marginal FLOPs
on \emph{more} thinking-token budget pay better than spending them on
\glsterm{rlvr-3}{RLVR} training, on parallel sampling (\cref{sec:self-consistency}), or
on a verifier-guided search (\cref{sec:tree-search})? s1's headline
suggests the trade-off is real --- 7 GPU-hours of SFT plus a longer
$L$ at inference beats 394 GPU-hours of SFT at the same inference
length \citep[\S 5.1]{s1-2025} --- but the comparison is one slice
through a many-dimensional surface. \Cref{sec:search-vs-rl} picks up
the search-vs-RL fixed-compute trade-off directly, transcribing
Snell's FLOPs-matched curves \citep[\S 6]{snell2024} for the
inference-vs-pretraining axis and lining them up against
\glsterm{dapo-2}{DAPO}'s step-efficiency and R1's training-token costs for the
inference-vs-RL axis. The thinking-token budget $L$ that this section
formalised becomes one of three FLOPs-bearing variables on that
trade-off surface, alongside the parallel-sample count $N$ of
\cref{sec:self-consistency} and the \glsterm{rlvr-3}{RLVR} step count of
\cref{sec:rlvr}.
